%%\section{Results}

%%%%%%%%%%%%%%%%%%%%%
\begin{figure}[!htpb]
\centering
\includegraphics[trim = 5mm 12mm 18mm 8mm, clip, width=0.48\columnwidth]{RatioResultPlot_CrossSectionRatio_Enu_z06_91_Nu_v19p1}
\includegraphics[trim = 1mm 12mm 18mm 8mm, clip, width=0.48\columnwidth]{RatioResultPlot_CrossSectionRatio_x_z06_91_Nu_v19p1}
\includegraphics[trim = 5mm 12mm 18mm 8mm, clip, width=0.48\columnwidth]{RatioResultPlot_CrossSectionRatio_Enu_z26_92_Nu_v19p1}
\includegraphics[trim = 1mm 12mm 18mm 8mm, clip, width=0.48\columnwidth]{RatioResultPlot_CrossSectionRatio_x_z26_92_Nu_v19p1}
\includegraphics[trim = 5mm 00mm 18mm 8mm, clip, width=0.48\columnwidth]{RatioResultPlot_CrossSectionRatio_Enu_z82_93_Nu_v19p1}
\includegraphics[trim = 1mm 00mm 18mm 8mm, clip, width=0.48\columnwidth]{RatioResultPlot_CrossSectionRatio_x_z82_93_Nu_v19p1}
\caption{Ratios of the charged-current inclusive \numu cross section per nucleon as a function of \Enu (left) and as a function of reconstructed \Xbj (right)
for C/CH (top), Fe/CH (middle), and Pb/CH (bottom).  Error bars on the data (simulation) show the statistical (systematic) uncertainties.  The $\chi^{2}$ calculation includes correlations among all bins shown.  Events with \Xbj greater than 1.5 are not shown.}
\label{fig:xsec_ratio_plots}
\end{figure}
%%%%%%%%%%%%%%%%%%%%%
Ratios of charged-current \numu cross sections per nucleon $\sigma\left(\Enu\right)$ and $\frac{d\sigma}{d\Xbj}$ are shown in Fig.~\ref{fig:xsec_ratio_plots}\SuppNote{See Supplemental Material\SuppLocation for cross section ratio measurements compared to simulation in tabular form and correlations of uncertainties among bins}.
Simulation reproduces these measurements within roughly 10\%.  
In contrast, measurements of $\frac{d\sigma^{\A}}{d\Xbj}/\frac{d\sigma^{CH}}{d\Xbj}$ show a suppression of the ratio compared to simulation at low \Xbj and an enhancement at high \Xbj, both of which increase with the size of the nucleus.

Low \Xbj bins are expected to show shadowing, which lowers the cross section for heavier nuclei~\cite{e665,Allport1989417,Kopeliovich:2012kw}.
Shadowing in these data may be larger than predicted for several reasons.
First, our data are at low \Qsq in the nonperturbative range (80\% of events below \unit[1.0]{\GeVcsq} and 60\% below \unit[0.5]{\GeVcsq}), while the model is tuned to data at much higher \Qsq where shadowing is well measured.
Second, shadowing in the model is assumed to be the same for C and Pb and equal to measurements from Fe~\cite{Bodek:2004pc}.
Finally, the shadowing model used for comparison is based on charged lepton data, which do not have axial-vector contributions.
The array of nuclear models available to modern neutrino experiments give 
similar results for these cross section ratios; none of which is confirmed by the data.

Higher \Xbj bins contain mostly ($\textgreater63\%$) quasielastic events, 
whose rates may be enhanced by meson-exchange currents~\cite{Carlson:2001mp,PhysRevC.80.065501,PhysRevC.81.045502,Bodek:2011ps,Amaro2011151,PhysRevC.83.045501,PhysRevD.88.113007},
which are not in the simulation.
The excess observed here may be related to the excess in MINOS Fe data at low inelasticity compared to a simulation with nuclear corrections based on lighter nuclei similar to GENIE's~\cite{PhysRevD.77.072002,arxiv:0708.1495}.
The failure of nuclear scaling models in this region has profound implications for neutrino oscillation experiments that utilize quasielastic events.
For example, T2K~\cite{PhysRevLett.111.211803,PhysRevLett.112.061802} must apply a nuclear model to relate the rate in the carbon of a near detector to oxygen in the far detector.  
LBNE~\cite{lbne_new} must extrapolate existing data on C, Fe, Pb to Ar.
Until better models exist that cover the relevant kinematic domain, oscillation experiments must 
incorporate the discrepancies measured here in evaluating systematic uncertainties.
More theoretical work is needed to correctly model nuclear effects in neutrino interactions, from the quasielastic to the deep inelastic regime.
